\chapter*{Prakata}
\addcontentsline{toc}{chapter}{Prakata}

Buku ini disusun dari satu kebiasaan: tidak ada angka yang ditulis sebelum
diukur. Tiap bab berangkat dari satu pertanyaan rekayasa, menjalankan
program yang menjawabnya di komputer penulis, lalu mencatat hasilnya apa
adanya, termasuk hasil yang melawan harapan.

Kebiasaan itu punya akibat yang terlihat di seluruh halaman. Angka pada
tabel bukan angka dari buku lain, melainkan keluaran program yang
disertakan. Kode pada setiap listing adalah salinan dari file yang
benar-benar dijalankan, bukan kode yang hanya hidup di halaman cetak. Ketika
nama di dalam kode berubah dan nama itu ikut tercetak pada keluarannya,
programnya dijalankan ulang dan angkanya diperbarui.

Isinya ditujukan kepada dua pembaca. Pembaca pertama adalah pengembang yang
sudah membuat aplikasi web dan ingin tahu mengapa aplikasinya lambat, mahal,
atau rapuh, serta bagian mana yang pantas diperbaiki lebih dahulu. Pembaca
kedua adalah pengajar dan mahasiswa yang membutuhkan bahan yang dapat
dijalankan ulang, bukan dibaca saja. Bagi pembaca kedua, seluruh prosedur
pengukurannya dikumpulkan sebagai lampiran di belakang, lengkap dengan
lembar isian dan pertanyaan yang jawabannya menuntut pembacaan angka.

Seluruh program, data uji, dan perintah yang dipakai di buku ini tersedia di
repositori yang menyertainya, tersusun per bab di dalam direktori
\berkas{sources/}. Bab Metode Pengukuran menjelaskan lingkungan tempat
angkanya diambil dan cara mengulanginya.

Terima kasih kepada mahasiswa Program Studi Informatika Universitas Pradita
yang menjalankan ulang pengukuran di buku ini pada komputer mereka
masing-masing. Angka mereka berbeda dari angka di sini, dan perbedaan itulah
yang membuat pola perbandingannya menjadi pelajaran, bukan hafalan.
